% difference between two pdf files: https://draftable.com/compare
% insert a pdf: https://tex.stackexchange.com/questions/8662/merge-two-pdf-files-output-by-latex
%   \usepackage{pdfpages}
%   \includepdf[pages=-]{ChangesFromV01toV02.pdf}
% gs -sDEVICE=pdfwrite -dCompatibilityLevel=1.4 -dNOPAUSE -dQUIET -dBATCH      -sOutputFile=ResponseCompressed.pdf Response.pdf
%
\documentclass[conference, onecolumn, 10pt]{../IEEEtran}
\IEEEoverridecommandlockouts
\usepackage{times}
 

\usepackage{microtype,xparse,tcolorbox}
\usepackage{mathtools, cuted}
\usepackage{todonotes}
\usepackage{graphicx}
\usepackage{wrapfig}
\usepackage{hyperref}
\usepackage{tikz}
\usepackage{pdfpages}
\definecolor{myMagenta}{RGB}{233,63,246}
\definecolor{myDarkGreen}{RGB}{97,173,89}
\providecommand{\norm}[1]{\left\lVert#1\right\rVert}
\pagestyle{plain}
\newcounter{reviewer}
\setcounter{reviewer}{0}
\newcounter{point}[reviewer]
\setcounter{point}{0}
\usepackage{lastpage}


\setlength\parindent{0pt}

\renewcommand{\thepoint}{[R\,\thereviewer.\arabic{point}] } 

\newcommand{\response}[1]{\refstepcounter{point}  \smallbreak \noindent
	\hangindent=3.5em
	\hangafter=1
	\thepoint \space ``\textcolor{blue}{#1}" \par\smallskip}

\newcommand{\letter}[2]{\vspace*{\fill} 
	
	\begin{quote}  
		
		\includegraphics[width=8cm]{engineeringLogo.eps}
			
		
	
		\begin{flushright}
		Aaron T.\ Becker\\Department of Electrical Engineering\\University of Houston\\Houston, TX 77004, USA\\Email: atbecker@central.uh.edu\\Web: youTube.com/aabecker5\\
		
		\today \vspace{.1cm} \\
		\end{flushright}
		
		#1\medskip\\ 
	
		Dear #2 Editorial Office,\medskip\\
		Please find attached the revised paper, #1, previously submitted as paper number 18-0277,
		along with the document containing a response to the reviewers. We are grateful to the reviewers for helping us improve our manuscript through their comments and questions. Please let us know if further information is required. This work extends the preliminary conference paper, ``\emph{Algorithms for shaping a particle swarm with a shared input by exploiting non-slip wall contacts}'', presented at the 2017 IEEE/RSJ International Conference on Intelligent Robots and Systems (IROS). The conference paper considered only square workspaces. This work extends the analysis to convex workspaces and 3D positioning. This paper also implements the algorithms using a hardware setup inspired by the anatomy of the gastrointestinal tract.\medskip
		\medskip
		
		The new revision of our paper includes multi-media so that others can build on our results. These include illustrative videos for the simulations and experiments: 
	\begin{itemize}	 
 \item Video  \emph{01Model.mov} animates the concepts of uniform control inputs and non-slip boundary contacts.
 \item Video  \emph{02DeltaConfigurationSpace.mp4} demonstrates how the $\Delta$ configuration space (the configuration space showing the difference in position for two particles) is constructed for a variety of workspaces.
 \item Video  \emph{03SimulationWorkSpaces.mp4} shows demonstrations of motion planning in square and disc-shaped workspaces. 
 \item Video  \emph{04Hardware Experiments.mp4} shows an experiment trial moving two particles to goal positions in a small intestine phantom, then an experiment trial using cow stomach tissue.
		\end{itemize}	 	
		 
		   We also include four  Mathematica Notebooks (.nb files) containing the simulation code.  This code generates plots for the paper and animates the path planner algorithm for arbitrary inputs.
		   	\begin{itemize}	 
 \item   \emph{SquareWorkSpace.nb} generates paths in a square workspace for two particles.
 \item  \emph{CircularWorkSpace.nb} generates paths in a circular workspace for two particles.
 \item   \emph{DeltaConfigurationSpacePolygon.nb} generates the $\Delta$ configuration space for convex polygons
  \item   \emph{ShortestPathForADisk.nb} analytical solution for shortest path that touches a boundary from one position to another position in a circular workspace.
		\end{itemize}	 
		   
		%   At the end of this letter, pages~\pageref{sec_lastpage}--\pageref{LastPage} is the revised manuscript, formatted so all changes can be easily identified by the reviewers, by using green colored text to indicate revised passages. 
		 The revised manuscript is formatted so all changes can be easily identified by the reviewers by using green colored text to indicate revised passages. 

		
		
		
		Sincerely,\\
		
	
			\includegraphics[width=4cm]{SignatureAaronTBeckerWhite.png}
			\medskip\\
		Aaron T.\ Becker (on behalf of all the authors) 
	\end{quote}
	\vspace*{\fill}
	\break }

\NewDocumentCommand \Reviewer { m } {
	\stepcounter{reviewer}
	\subsection*{Comments~by~Reviewer~#1}
}

\begin{document}
	\letter{\emph{Exploiting Non-Slip Wall Contacts  to Position Two Particles Using The Same Control Input}}{IEEE Transactions on Robotics}
	
	\section*{Response to Reviewers}
	In the following document, we have provided detailed responses to the comments and questions of the reviewers. Comments and questions by reviewers are in \textcolor{blue}{blue}, our responses are in black.
	
	
	
	\emph{Comments by Associate Editor}
	\response{The Associate Editor (AE) was able to obtain two high-quality reviews
from researchers working directly in this field. After considering
their reviews, it is the opinion of the AE that this paper is not
currently suitable for publication in T-RO. However, the authors are
encouraged to prepare a revised draft, addressing the concerns of the
reviewers.  Please also address how the technique could be practically applied to
navigate inside an actual 3D intestine as opposed to a cross section in
which the particles move on a planar surface. Would magnetic forces be
sufficient to overcome gravity? Would the particles need to be made
neutrally buoyant?}

	We are grateful for the efforts of our reviewers. We have thoroughly revised this paper by (1) restructuring the paper as suggested by the reviewers (2) constructing reachability sets for all initial conditions, proving controllability as long as the initial positions are distinct and the final positions are not   antipodes, and (3) updating almost all the figures and captions. 
	
	We made the model for both 2D and 3D more clear -- the particles are assumed to be neutrally buoyant so that the same force is applied to each particle.
	Please see the detailed responses below.
	
	
	
	\Reviewer{\#1}\response{This paper explains an underactuated control strategy for small teams
of magnetic microparticles which move under the influence of a single
broadcast signal. This underactuated control problem is relevant to
microrobotic swarms moving under magnetic field input. In previous
work, the authors have explored interesting control methods relying on
obstacles in the environment, but the case studied here is more
application-realistic. The control method here is clever and
successful. The work is interesting to read and I think that anyone
with a basic background in controls will find this work interesting.  
The main concept used here to achieve multi-particle control is that
the particles experience a no-slip condition when pushed against the
environment wall. The particles can be removed from the wall by pulling
them. 
The paper is an extension of an IROS2017 paper which introduced the
idea for square workspace only. This TRO submission generalizes to a
non-square workspace, which is a notable improvement. This version of
the work will thus be a helpful addition to the literature. 
The paper gives a very nice review of previous work in using broadcast
signals for multi-particle control.  }
	Thank you! We have completed a major revision, as explained below.
	
	\response{Is it expected that the no-slip condition will be a strongly satisfied
assumption for in-vivo biomedical applications? In cases where there
were some slip, how could that be handled?}
	
	To maintain clarity, this paper focused on systems with non-slip boundaries.  While several environments with non-slip boundaries exist (e.g. tripe, and intestines with villi), extending to lower friction values is an intriguing avenue that unfortunately does not fit within the 12 page limit for this journal, but is something we are investigating. The reachable sets with non-infinite friction are smaller. 
The conclusion now explains `` This paper assumed friction was sufficient to completely stop particles in contact with the boundary. The algorithms would require retooling to handle small friction coefficients."
	
	\response{This control method requires advance knowledge of the workspace
boundary geometry, in addition to micro-particle state feedback. It
would be interesting to hear about how accurate the boundary geometry
must be known for this control to work well, and perhaps compare this
will medical imaging techniques such as CT or MRI which could
potentially be used to generate such maps. For motion through 3D
lumens, the cross-sectional geometry will change throughout the lumen.}
	
	These are excellent points. This paper assumes perfect knowledge of the boundary and initial conditions of the particles.  We added text explaining that systems with large numbers of protrusions or concave workspaces are best handled by motion planners such as RRT.
	
	Section IV-F: ``Workspaces that are tortuous or with many obstacles are better handled by other path
planners, such as RRT [29], or [5], which used collisions with protrusions of the workspace to rearrange particles.''.	
	
	Also, during the hardware experiments, the particle that was pushed into the com stomach boundary disappeared from the camera, but control still was successful, indicating that perfect sensing is not necessary: ``Our algorithm successfully delivered the particles to goal positions in 5 out of 5 trials.''

	\response{Abstract: It is stated that ``given 3 orthogonal magnetic fields''.
This statement should be made more precise, because the magnetic
pulling here uses field gradients, not fields. In addition, these
gradients are not orthogonal. An accurate statement would need to be
more complex.}
	
	Yes.  To make this work as broadly applicable as possible (magnetic inputs are just one example) we changed the abstract to say ``This paper investigates particle control with uniform forces (the same force is applied everywhere in the workspace). Given a controllable field that can generate bidirectional forces in three orthogonal directions, steering one particle in 3D is trivial. ''
	  
	  	
	\response{Fig 2: ``first contact point'' is not clear. There is no way given to
tell where these line up with actual positions on the workspace. }
	
	  Fig 2 is now Fig. 11. We rearranged the order because the optimization results can be applied to Alg.~1, but are not necessary to implement or understand the algorithm.
	    We updated this figure to clarify what the contact points are by adding and labelling tick marks from 0 to 4 along the boundary of the square workspace and the distance plot.  We also added angle markings and green points showing the optimal solution in both the workspace and the distance plots. 
	
	
	\response{Fig 3: the meaning of grey areas should be mentioned in the fig or
caption.}
This is now Figure 2.  We added a legend and caption text stating:
``Gray areas denote regions unaccessible by our motion planner. The particle start positions must be distinct $(\norm{s_2-s_1} \geq \epsilon)$, and at least one goal position must be farther than  $\epsilon$  from the boundary, where $\epsilon$ is a small but nonzero user-specified constant.''

	
	\response{Figures are hard to read because there is a lot going on, the font is
often too small and text sometimes is obscured.}

There were a number of problems with the initial figures.  We have added legends to Figs. 2, 3, 5, 15, and added captions with visual elements to 2, 14, 16.  We standardized our variable names and icons, and added Figs~8 and 9.

	
	\response{III.C. It�s not clear why this section is relevant. Why is the shortest
path which intersects the wall needed? Some writing structure could
help the reader understand what the goal is here.} 

	To better present the motion planner, we moved this section, which covers an optimization result, to the end, and altered the framing text:
	``Algorithm 1 provided a technique to bring two particles to goal positions using global inputs, but did not optimize path length. Changing the relative positions of particles in any workspace requires making one particle contact the boundary. In this section we present two results that can be incorporated into Algorithms 2 and 3 to generate shorter motion paths.''
	
	

	
	
	\Reviewer{\#6}
	\response{In this manuscript algorithms to position two particles in arbitrary
locations under uniform actuation are proposed.  The algorithms rely on
differentiating between particles by contacting them with walls or
boundaries of the domain, at which they are assumed to not move unless
actuated in a direction away from the wall.  The paper presents some
results on shortest paths that contact surfaces (but see below on some
presentation issues), then describes the algorithm, which basically
entails moving both particles so that one particle contacts a boundary,
adjusting the relative displacement between particles while the one
particle remains at the boundary, then translating both particles to
the desired location.
Simulations of the algorithm are presented in square and circular
domains, and experiments are described in circular cross sections
inspired by intestinal and stomach scenarios.
%
Although similar ideas were presented earlier by the authors using
either an obstacle in the workspace or square workspaces, this paper
extends that work to circular and convex polygonal workspaces.	The
overall idea is a novel contribution as other swarm control techniques
focus on either heterogeneity of microrobots or of actuation (such as
fields) rather than distinguishing particles by their proximity to
boundaries.  However, in my opinion the manuscript requires major
revisions before being suitable for publication, mainly to address the
theoretical effectiveness of the algorithm (as detailed below), and
less so to address clarity of presentation. }
	Thank you. We were pleased to get this review, because it led us to think more deeply on how to present the results and led to the major revisions detailed below.
	
	\response{  The main weakness of the paper in my opinion is that the authors
never really prove that their algorithm works, or alternatively how
general their algorithm is, i.e. what are its limits of applicability. 
To be more specific, in their Algorithm 1, if the desired configuration
is not within the 2-move reachable set, then the algorithm targets
instead the closest point in the 2-move reachable set, and then
"iterates until we reach the goal."  When does such iteration actually
achieve the goal?  What are the attainable final goals that can be
reached after iteration?  Does the space of attainable goals depend on
the initial positions of the particles?  While there are simulations,
the ones presented do not explore the entire possible space (which is
understandable since they it is quite large).  I think an analysis of
the attainable space of their algorithm is needed.}
	
	An excellent suggestion!  This was a major oversight in the first version.  It is  possible to explore the possible space, and we can be quite clear about the locations that are attainable.
	We rearranged the paper and added Section IV-F, which analyzes the reachable set:
	
	``The $\Delta$ configuration enables an iterative method to compute the accessible workspace.  
 Due to symmetry of the workspace, the fraction of the $\Delta$ configuration space reachable in $2 k$ moves is a function of only the initial separation distance $d_{12}$.
 The angle $\theta$ between the initial particle positions simply rotates the reachable $\Delta$ configuration space. As long as the initial configurations are distinct $(s_1 \neq s_2)$, the reachable set grows quickly. 
  This relationship is shown in Fig.~15.
   Only antipodal locations are unreachable $(\norm{g_2-g_1} = 1)$, but can be asymptotically approached.
   Indeed, even with a tiny initial separation of $d_{12} = 0.001$, after 14 moves 90\% of the $\Delta$ configuration space is reachable. 
   In two moves, the maximum reachable fraction of 0.373 is achieved with $d_{12} \approx 0.81$." 
   
   We added Figs. 8 and 9 which illustrate the effects of 2-move reachable sets, as well as 4-move, 6, 8, 10 , 12, and 14-moves.
	
	Fig. 15 analyzes the worst case: when the goal locations are at the antipodes.  Even in this situation we can approach arbitrarily close to the desired configuration, but can never achieve it. This figure shows that the approximation error decrease can be fit by a 1/(distance travelled$^3$).
	 
	
	\response{There are issues with the presentation that can be easily improved.
 At the beginning of the paper many terms are not defined which makes
figures and discussion hard to follow. Definitions do come later, but
they should be moved up.  For example, in Fig 2 which symbols are
targets and which are initial conditions are not defined.  In Fig 3
epsilon is shown but not defined nor is its significance explained.  s
and g and $\Delta s$ and $\Delta g$ are not defined. In Fig 12, where is
the goal on the boundary? Can that be indicated on the figures?}

Thank you for detailed points. We performed a large restructuring. 
Fig 2 is now 11. We updated this figure to clarify what the contact points are by adding and labelling tick marks from 0 to 4 along the boundary of the square workspace and the distance plot.  We also added angle markings and green points showing the optimal solution in both the workspace and the distance plots. We also added visual icons to the caption: ``from starting positions 
 (\protect\tikz \protect\draw[red,fill=white,line width=0.3mm] (-1ex,0) rectangle (0,1ex);,
 \protect\tikz \protect\draw[blue,fill=white,line width=0.3mm] (-1ex,0) rectangle (0,1ex);)
 to goal positions  
 (\protect\tikz \protect\draw[red,fill=white,line width=0.3mm] (0,0) circle (.5ex);,
 \protect\tikz \protect\draw[blue,fill=white,line width=0.3mm] (0,0) circle (.5ex);)''
 
 
 Figure 3 is now 2, we added both a legend to the figure and a descriptive caption: 
 ``Particles move from start positions 
(\protect\tikz \protect\draw[blue,fill=white,line width=0.3mm] (-1ex,0) rectangle (0,1ex);,
\protect\tikz \protect\draw[myMagenta,fill=white,line width=0.3mm] (-1ex,0) rectangle (0,1ex);), to goal positions (\protect\tikz \protect\draw[blue,fill=white,line width=0.3mm] (0,0) circle (.5ex);,
\protect\tikz \protect\draw[myMagenta,fill=white,line width=0.3mm] (0,0) circle (.5ex);).  Dashed lines show the shortest route if particles could be controlled independently.  Solid arrows show path given by  Alg.~1. Gray areas denote regions unaccessible by our motion planner. The particle start positions must be distinct $(\norm{s_2-s_1} \geq \epsilon)$, and at least one goal position must be farther than  $\epsilon$  from the boundary, where $\epsilon$ is a small but nonzero user-specified constant.''


 Figure 12 had the wrong caption.  It now is Fig. 14, and has the caption: ``Contour plots showing the number of moves and distance commanded if red particle's goal position is varied in $x$ and $y$. 
Starting positions of red and blue particles  
 (\protect\tikz \protect\draw[red,fill=white,line width=0.3mm] (-1ex,0) rectangle (0,1ex);,
 \protect\tikz \protect\draw[blue,fill=white,line width=0.3mm] (-1ex,0) rectangle (0,1ex);)
 and goal position of blue particle  \protect\tikz \protect\draw[blue,fill=white,line width=0.3mm] (0,0) circle (.5ex); are fixed.
The top row has the blue particle's goal position  at the origin, generating symmetric contour plots.
Moving the  blue particles' goal position  to $(-0.2,0)$, generates non-symmetric contour plots.''

We also added captions for start and end (colored squares and circles) to each drawing.

	
	\response{(more presentation) At the beginning of section IV it wasn't clear
to me whether Algorithm 1 had been described yet when it was first
mentioned, and whether that sentence was a description of what was to
come, or was as assertion that logically followed from the previous
parts of the paper.  }
	
	We  rewrote this section. It now begins with:
	``This section presents an algorithm, Alg.~1, that uses non-slip contacts with walls to arbitrarily position two particles in a convex workspace. 
 Workspaces are 2D convex polygons with no internal obstacles. ''
		
	\response{A minor comment, sections IIIB and C are a little confusing since it
is not motivated why shortest path is being considered in the overall
argument of the paper.	Furthermore, it seems that shortest path in
IIIB is used for the situation moving two particles, while in IIIC it
is used for the situation of moving only one particle.	Consistency in
the presentation would help..}
	
	We rearranged the order because the optimization results can be applied to Alg. 1, but are not necessary to implement or understand the algorithm.  The results were combined into a new Section V, which opens by stating:
	
	``Algorithm 1 provided a technique to bring two particles to goal positions using global inputs, but did not optimize path length.   Changing the relative positions of particles in any workspace requires making one particle contact the boundary.
  In this section we present two results that can be incorporated into Algorithms 2 and 3 to generate shorter motion paths.''


	
%\label{sec_lastpage}\includepdf[pages=-]{ChangesFromV01toV02.pdf}	
	
	
		
\end{document}